\section{Introduction}

Quantum computation, sensing, and materials design fundamentally depend on the controlled use of quantum resources, including entanglement, coherence, superposition, and interference~\cite{NielsenChuang,Horodecki2009,Streltsov2017}. In most current approaches, these resources are treated as quantities to be maximized, leading naturally to strategies that explore increasingly large Hilbert spaces~\cite{Preskill2018,Arute2019}. While powerful in principle, this paradigm introduces severe scaling challenges and often obscures the physical structure underlying useful quantum dynamics~\cite{Eisert2010,Altman2015}.

An alternative perspective is that quantum advantage arises not from the size of Hilbert space, but from the \emph{organization} of quantum resources within it~\cite{Amico2008,Swingle2018}. In this view, the central problem is not to generate maximal entanglement, but to engineer entanglement that is structured, localized, and aligned with the geometry of the task~\cite{Verstraete2008,Orus2014}. This shift reframes quantum control as a problem of resource organization rather than resource amplification. Useful quantum dynamics are then those that remain confined to structured manifolds where entanglement, interference, and correlation are jointly optimized~\cite{KimHuse2013,Choi2016}.

Here we introduce Entanglement-Dynamic Topology (EDT) as a unifying framework for programmable quantum resource engineering, building on recent developments in curvature-driven and topologically structured quantum dynamics~\cite{Hu2024_ACP,Hu2025_APQ,Hu2026_JPCM}. EDT describes a regime of quantum dynamics in which resources are simultaneously \emph{structured, interacting, topological, and dynamical}. Structure, arising from symmetry, constrains quantum states into reduced manifolds~\cite{Wigner1959,Georgi1999}. Interaction, driven by non-adiabatic coupling such as electron--phonon interaction, generates and redistributes entanglement~\cite{Mead1979,Yarkony1996}. Topology organizes this structure through nodes, poles, and vortex-like configurations in parameter space~\cite{Berry1984,Xiao2010}. Dynamics, through modulation of band energies and electronic potentials, renders these manifolds controllable and programmable~\cite{Bukov2015,Eckardt2017}.

A key feature of EDT is that it defines a finite regime rather than a fine-tuned point. Within this regime, quantum dynamics exhibit coherent, structured behavior, while outside it the dynamics become either weakly correlated or overdriven and incoherent~\cite{Larkin1969,Swingle2018}. This regime-based perspective provides a robust and experimentally meaningful foundation for quantum resource engineering, in contrast to approaches that rely on asymptotic scaling or finely tuned optimality.

To connect this framework to real physical systems, we first develop a data-driven method for identifying EDT directly from electronic-structure observables. Using K$_3$C$_{60}$ as the main material example, we analyze curvature-resolved band structure and show that EDT emerges as a finite momentum interval in which paired curvature branches admit a consistent local quadratic description. Supplementary archived comparison datasets for V, Na, LaH$_{10}$, and YBa$_2$Cu$_3$O$_7$ are then used as portability checks of the same diagnostic rather than as co-equal material claims. This provides an operational criterion for detecting EDT and establishes a direct link between microscopic electronic structure and structured quantum resources~\cite{Gunnarsson1997,Capone2009}.

We then show that this regime arises generically in a minimal dynamical model, where loop-induced operator spreading remains structured and compatible with persistent many-body correlation, and that it remains robust under parameter variation, bounded circuit complexity, and approximate implementation~\cite{Lloyd1996,Childs2019}. Finally, we demonstrate the same finite-regime logic in NV-center XY8 sensing, where optimal performance emerges from balancing filter response, coherence, pulse fidelity, and readout overhead within a bounded control window~\cite{Degen2017,Barry2020,Casola2018}.
